\documentclass[11pt]{amsart}
\usepackage{amsmath,amssymb,amsthm,mathtools,mathrsfs}
\usepackage[margin=1.08in]{geometry}
\usepackage{booktabs}
\usepackage{array}
\usepackage{tabularx}
\usepackage[hypertexnames=false]{hyperref}
\hypersetup{hidelinks}
\usepackage[final,expansion=false]{microtype}
\usepackage{url}
\setlength{\parskip}{4pt plus 1pt}
\emergencystretch=2.5em

\newtheorem{theorem}{Theorem}[section]
\newtheorem{proposition}[theorem]{Proposition}
\newtheorem{lemma}[theorem]{Lemma}
\newtheorem{corollary}[theorem]{Corollary}
\newtheorem{mainthm}{Theorem}
\renewcommand{\themainthm}{\Alph{mainthm}}
\theoremstyle{definition}
\newtheorem{definition}[theorem]{Definition}
\theoremstyle{remark}
\newtheorem{remark}[theorem]{Remark}

\newcommand{\C}{\mathbb C}
\newcommand{\R}{\mathbb R}
\newcommand{\N}{\mathbb N}
\newcommand{\Z}{\mathbb Z}
\newcommand{\Q}{\mathbb Q}
\newcommand{\F}{\mathbb F}
\newcommand{\Rel}{\operatorname{Re}}
\newcommand{\Iml}{\operatorname{Im}}
\newcommand{\supp}{\operatorname{supp}}
\newcommand{\DW}{\mathcal D^{W}}
\newcommand{\PW}{\mathcal P^{W}}
\newcommand{\Nr}{\mathcal N}
\newcommand{\Ef}{\mathcal E}
\newcommand{\I}{\mathcal I}
\newcommand{\dg}{\deg}
\newcommand{\Kc}{\kappa_C}
\newcolumntype{Y}{>{\raggedright\arraybackslash}X}

\begin{document}

\title[An intersection frame for the explicit formula]
{An Intersection Frame for the Explicit Formula:\\
Hodge-Index Transfer, a Frame Defect Theorem\\
and a Function-Field Model}

\author{David Alejandro Trejo Pizzo}
\thanks{Independent Researcher. \texttt{dtrejopizzo@gmail.com}.}

\date{\today}

\begin{abstract}
Let $\DW$ be the space of real $C^2$ functions on $\R$ which, with their first two derivatives, decay
like $e^{-(1+\varepsilon)|t|}$, and let $\I$ be the symmetric bilinear form on $\DW$ defined by
Weil's explicit formula, so that $\I(F,G)=\Phi_F(0)\Phi_G(1)+\Phi_F(1)\Phi_G(0)-\sum_\rho\Phi_F(\rho)\Phi_G(1-\rho)$
with $\Phi_F(s)=\int F(t)e^{(s-1/2)t}\,dt$. We equip $(\DW,\I)$ with the structure Weil's proof of the
Riemann hypothesis for curves uses on $C\times C$: a degree, an effective cone (pointwise nonnegative
functions), a radical, an ample class and an affine duality. The class $K$ of Riemann's integral
representation in P\'olya's normalization ($\Phi_K=\xi$) is a degree-computing ample class, the
radical is cut out by the polar and zero data, and $\DW$ splits orthogonally into a hyperbolic plane
and the primitive sector. Hartshorne's proof of the Hodge index theorem transfers verbatim:
a function $h^0$ with a positivity clause, an effectivity clause and a Riemann--Roch inequality would
give $\I(F,F)\le0$ on every primitive compactly supported test. We prove that no such function exists.
The obstruction is a cosh inequality on the effective cone, with sharp constant
$0.99846\ldots<1$ at the first zero of $\xi$, which produces explicit classes (of degree $\ge325$ when $\chi=0$) of positive self-intersection with
no effective representative; beyond classical facts, the proof uses only the numerically known
simplicity and isolation of the first zero and Trudgian's explicit bound for $S(T)$. The same theorem holds, with defect $1-\sqrt{7/11}$ and a
witness in degree $3$, for an elliptic curve over $\F_2$, where every object of the frame exists and
the Hodge index theorem is true. In the curve model the frame's effective cone is exactly the cone
on the graphs of the Frobenius powers, the class $K$ is half the ample ruling class, and the missing
effective classes are the graphs of sums $m+n\pi$ and the fibres.
\end{abstract}

\maketitle
\tableofcontents

%======================================================================
\section{Introduction}\label{sec:intro}

Weil's proof of the Riemann hypothesis for a curve $C$ over a finite field runs on the surface
$C\times C$: Riemann--Roch and Serre duality make every class of positive self-intersection and
positive degree big, and the Hodge index theorem follows \cite[Ch.~V, 1.9]{Hartshorne}. Weil's
explicit formula over $\Q$ supplies a quadratic form with the same polar shape. This paper builds
around that form the minimal intersection-theoretic structure the surface proof consumes, proves
what can be proved about it, and shows exactly where it fails. The failure is a property of the
frame, not of the zeros: it occurs verbatim on a curve over $\F_q$.

\subsection{The frame}
Write $\Phi_F(s)=\int_\R F(t)e^{(s-\frac12)t}\,dt$ and let $\rho$ run over the zeros of
$\xi(s)=\frac12s(s-1)\pi^{-s/2}\Gamma(s/2)\zeta(s)$ with multiplicity. On the Weil class $\DW$
(Definition~\ref{def:DW}) the explicit formula defines a symmetric bilinear form $\I$ with
\[
\I(F,G)=\Phi_F(0)\Phi_G(1)+\Phi_F(1)\Phi_G(0)-\sum_\rho\Phi_F(\rho)\Phi_G(1-\rho),
\]
the zero sum absolutely convergent (Theorem~\ref{thm:comparison}); the definition itself is by the
arithmetic side and names no zero. The degree is $\dg F=\Phi_F(0)+\Phi_F(1)=\int F(t)\,2\cosh(t/2)\,dt$,
a class is \emph{primitive} if $\Phi_F(0)=\Phi_F(1)=0$, \emph{effective} if $F\ge0$ pointwise, and
$\Nr$ denotes the radical of $\I$.

\subsection{Main results}
\begin{mainthm}[structure; Sections~\ref{sec:polya}--\ref{sec:degree}]\label{thmA}
Let $K(t)=\sum_{n\ge1}(4\pi^2n^4e^{9t/2}-6\pi n^2e^{5t/2})e^{-\pi n^2e^{2t}}$. Then $K\in\DW$, $K>0$,
$\Phi_K=\xi$, $\dg K=1$, $\I(F,2K)=\dg F$ for all $F$, and $\I(K,K)=\frac12$. The radical is
$\Nr=\{N:\Phi_N(0)=\Phi_N(1)=0,\ \Phi_N(\rho)=0\ \forall\rho\}$; it is strictly larger than the ideal
$\Pi K\star\DW$ and meets $C_c^\infty(\R)$ only in $0$. With $K_{\rm odd}=(2\sinh\frac a2)^{-1}(\tau_aK-\tau_{-a}K)$,
$\DW=\mathrm{span}(K,K_{\rm odd})\oplus\PW$ is an $\I$-orthogonal sum of a hyperbolic plane and the
primitive sector. For every $F$ there is $C_F<\infty$ with $|\I(E,F)|\le C_F\dg E$ for all effective
$E$; hence $2K$ and every $F+2mK$, $m>C_F$, are ample.
\end{mainthm}

\begin{mainthm}[Hodge-index transfer; Theorem~\ref{thm:transfer}]\label{thmB}
Suppose $h^0:\DW\to\R_{\ge0}$ satisfies \textup{(b$'$)} if $h^0(E)>0$ then $E\equiv E'\bmod\Nr$
for some effective $E'\not\equiv0$, and \textup{(c)} $h^0(E)+h^0(K-E)\ge\frac12\I(E,E)-\frac12\I(E,K)+\chi$
for a constant $\chi$. Then $\I(F,F)\le0$ for every real primitive $F\in C_c^\infty(\R)$.
\end{mainthm}

\begin{mainthm}[frame defect; Theorems~\ref{thm:defect}, \ref{cor:smalldomain}]\label{thmC}
Given that the first zero $\rho_1=\frac12+i\,14.1347\ldots$ of $\xi$ is simple and the only zero
with $|\Iml\rho|\le21$ in the upper half-plane, no function $h^0$ as in Theorem~\ref{thmB} exists,
for any $\chi$. The same holds for $h^0$ defined only on $C_c^\infty(\R)\oplus\R K$, or on
$\mathcal P^\infty_c\oplus\R K$. The witness is explicit: for
$F_0(u)=e^{-u^2/18}\sin(\gamma_1u)$, the class $E=53.119\ldots F_0+400K$ has $\dg E=400$,
$\I(E,E)=220>\frac12\dg E$, and no effective representative modulo $\Nr$.
\end{mainthm}

The mechanism is the cosh inequality (Lemma~\ref{lem:cosh}): for effective $E'$ and every $\gamma,\theta$,
$|\int E'(u)\cos(\gamma u+\theta)\,du|\le\frac12s(\gamma,\theta)\dg E'$ with
$s(\gamma,\theta)=\sup_u|\cos(\gamma u+\theta)|/\cosh(u/2)$, and $s(\gamma_1,\pi/2)=0.998460\ldots<1$,
whereas the Riemann--Roch inequality (c), through the affine symmetry $E\mapsto K-E$, forces
$h^0(E)>0$ on all classes of degree $>1$ with $\frac12\I(E,E)-\frac14\dg E+\chi>0$. The effective cone
is a moment cone and is strictly thinner than the positive cone.

\begin{mainthm}[the curve model; Section~\ref{sec:curve}]\label{thmD}
Let $C/\F_q$ be a curve of any genus $g\ge1$ with distinct Weil numbers. The frame transcribes to functions
on the lattice $\Z$; the form becomes the intersection form on $\sum_n\R\,q^{-|n|/2}\Gamma_{q^n}\subseteq
NS(C\times C)_\R$, with Gram matrix $q^{-|a-b|/2}N_{|a-b|}$ given by the point counts; the frame's
effective cone is the cone on the normalized Frobenius graphs and their transposes; the class with
$\Phi=\xi_C$, normalized to degree $1$, is $\frac12(F_1+F_2)$, half the ample ruling class. The frame
defect theorem holds on $C$. For $E/\F_2:y^2+xy=x^3+1$ the sharp constant is $\sqrt{7/11}$ and the
witness has degree $3$ (at $\chi=0$); the section cone has constant $1$, and the ruling $F_1$ is effective but not
in the frame's cone.
\end{mainthm}

Section~\ref{sec:xia} treats $\Xi_a(s)=\xi(s+a)\xi(s-a)$, $0<a<\frac12$, which has the same frame
and zeros off the line: there no $h^0$ exists either, for a different reason
(Theorem~\ref{thm:xia-noh}). Together with Theorem~\ref{thmD} this locates the defect: it is present
whether the Hodge index inequality is false ($\Xi_a$), true ($C/\F_q$) or open ($\xi$).

\subsection{Relation to other work}
The frame follows Weil's proof for curves \cite{Weil48} in Hartshorne's formulation \cite{Hartshorne};
the explicit formula is Weil's \cite{Weil52}, in the form of \cite{Barner1981} on larger classes.
The comparison identity on compactly supported tests is \cite[Thm.~4.1]{Paper45}, and the
construction of primitive tests from an off-line zero is \cite[Thm.~6.2]{Paper45}. This paper is
independent of the preceding paper \cite{Paper46} of the series. That paper studies the truncated Weil form on a finite window and
the positivity of its bottom on the primitive subspace, by dual certificates and finite reductions;
the present paper studies the extended form on the non-compactly supported class $\DW$ and asks
which intersection-theoretic structure it carries. The two meet only in the comparison identity:
on primitive compactly supported tests $\I(F,F)=-\sum_\rho\Phi_F(\rho)\Phi_F(1-\rho)$, which is the
form whose window positivity \cite{Paper46} studies. No result of either paper is used in the other,
and nothing here bears on the positivity problem of \cite{Paper46}.

\subsection{Conventions}
Numerical constants are computed in 40-digit arithmetic (\texttt{mpmath}) and are stated with their
margins; this is floating-point computation, not interval arithmetic. Classical inputs:
\begin{itemize}
\item[(E1)] $\zeta'/\zeta(s)=-\sum_n\Lambda(n)n^{-s}$, absolutely convergent for $\Rel s>1$.
\item[(E2)] $\xi$ is entire, $\xi(s)=\xi(1-s)$, $\xi(0)=\xi(1)=\frac12$, its zeros are the nontrivial
zeros of $\zeta$ with multiplicity, $0<\Rel\rho<1$, and $\zeta(s)\ne0$ for $\Rel s\ge1$.
\item[(E3)] $\xi$ has order $1$ with its Hadamard product; $N(T+1)-N(T)=O(\log T)$; $\xi'/\xi(\sigma\pm iX_j)=O(\log^2X_j)$
on suitable heights $X_j\in[j,j+1]$; $N(T)=O(T\log T)$, so $\sum_\rho|\gamma_\rho|^{-2}<\infty$
\cite[Ch.~15--17]{Davenport2000}.
\item[(E4)] Stirling's formula for $\psi$ in vertical strips.
\item[(E5)] Fourier inversion; (E6) the Paley--Wiener theorem.
\item[(X1)] P\'olya's form of Riemann's integral: $\xi(\frac12+i\tau)=\int K(t)e^{i\tau t}dt$ with $K$
even \cite[\S10.1]{Titchmarsh1986}, \cite{Polya1926}.
\item[(X2)] Jensen's formula: an entire function of exponential type has $O(r)$ zeros in $|s|\le r$.
\item[(X3)] Riemann--von Mangoldt asymptotics and a positive proportion of simple zeros
\cite{Conrey1989} (the proportion $\frac13$ of \cite{Levinson1974}, whose zeros are simple by \cite{HeathBrown1979}, suffices).
\item[(X4)] The first zero: $\rho_1=\frac12+i\gamma_1$, $\gamma_1=14.134725141734693790\ldots$, is simple,
and $\xi$ has no other zero with $0<\Iml\rho\le21$ (classical; see \cite[\S15]{Titchmarsh1986} and
\cite{PlattTrudgian2021}); and Trudgian's explicit
bound for $S(T)$ \cite{Trudgian2014}, which gives $\#\{\gamma\in[T,T+1]\}\le10\log T$ for $T\ge21$.
\item[(X5)] $\sum_\rho\Rel(1/\rho)=1+\frac12\gamma_E-\frac12\log4\pi$ \cite[Ch.~12]{Davenport2000}.
\item[(X6)] Hardy's theorem: $\xi$ has infinitely many zeros on $\Rel s=\frac12$ \cite{Hardy1914}.
\item[(X7)] For curves over $\F_q$: rationality, functional equation and the Riemann hypothesis for
$Z_C$; the Lefschetz relations $\Gamma_{q^m}\cdot\Delta=N_m$,
$\Gamma_{q^m}\cdot\Gamma_{q^n}=q^nN_{m-n}$, $\Gamma_{q^m}^t\cdot\Gamma_{q^n}=N_{m+n}$ ($m\ge n\ge0$);
Tate's theorem on $NS(E\times E)$ \cite{Weil48,Tate1966}.
\item[(X8)] Intersection theory, Riemann--Roch, Serre duality, the Hodge index theorem and the
bigness criterion on a smooth projective surface \cite[Ch.~V]{Hartshorne}, \cite[\S2.2]{Lazarsfeld}.
\end{itemize}

%======================================================================
\section{The Weil class and the extended form}\label{sec:form}

\begin{definition}\label{def:DW}
$\DW$ is the space of real $g\in C^2(\R)$ for which there is $\varepsilon>0$ with
$g^{(j)}(t)=O(e^{-(1+\varepsilon)|t|})$ for $j=0,1,2$. Put $g^\vee(t)=g(-t)$, $\tau_ag=g(\cdot-a)$,
$\star$ for convolution and $\Pi=-\partial_t^2+\frac14$.
\end{definition}

\begin{lemma}\label{lem:DW}
(i) $\DW$ is a commutative real $\star$-algebra, stable under $\vee$ and every $\tau_a$, and contains
$C_c^\infty(\R)$. (ii) For $g\in\DW$, $\Phi_g$ is holomorphic on
$S_\varepsilon=\{-\frac12-\varepsilon<\Rel s<\frac32+\varepsilon\}$ and
\begin{equation}\label{eq:decay}
|\Phi_g(s)|\le\frac{C_2(g,\sigma)}{|s-\frac12|^2},\qquad C_2(g,\sigma)=\int|g''(t)|e^{(\sigma-\frac12)t}dt,
\end{equation}
uniformly on closed substrips. (iii) $\Phi_{g^\vee}(s)=\Phi_g(1-s)$, $\Phi_{\tau_ag}(s)=e^{(s-\frac12)a}\Phi_g(s)$,
$\Phi_{g\star h}=\Phi_g\Phi_h$, and $\Phi_{\Pi g}(s)=s(1-s)\Phi_g(s)$ when $g^{(3)},g^{(4)}$ also obey the decay.
(iv) $\ker\Pi\cap\DW=0$.
\end{lemma}
\begin{proof}
(i) If $|g|,|h|\le Ce^{-(1+\varepsilon)|t|}$ then $|g\star h|(t)\le C^2(1+|t|)e^{-(1+\varepsilon)|t|}$, and the
derivatives fall on one factor. (ii) Absolute convergence and Morera; two integrations by parts, whose
boundary terms vanish by the decay for $j=0,1$, give $\Phi_g(s)=(s-\frac12)^{-2}\int g''(t)e^{(s-\frac12)t}dt$.
(iii) Substitutions, Fubini and integration by parts. (iv) $\ker\Pi=\mathrm{span}\{e^{\pm t/2}\}$.
\end{proof}

\begin{theorem}[explicit formula on $\DW$]\label{thm:EF}
Let $h\in\DW$ be real and even. Then $\sum_\rho\Phi_h(\rho)$ converges absolutely and
\begin{equation}\label{eq:EF}
\sum_\rho\Phi_h(\rho)=\Phi_h(0)+\Phi_h(1)+\frac1{2\pi}\int_\R\hat h(\tau)\Big[\Rel\psi\big(\tfrac14+\tfrac{i\tau}2\big)-\log\pi\Big]d\tau
-\sum_{n\ge2}\frac{\Lambda(n)}{\sqrt n}\big[h(\log n)+h(-\log n)\big],
\end{equation}
with $\hat h(\tau)=\Phi_h(\frac12+i\tau)$ and the prime sum absolutely convergent.
\end{theorem}
\begin{proof}
This is the contour proof of \cite[Thm.~3.3]{Paper45}, written there for $C_c^\infty$ tests, checked
against Definition~\ref{def:DW}. Integrate $\Phi_h\,\xi'/\xi$ over rectangles with vertical sides
$\Rel s=\frac32$ and $\Rel s=-\frac12$, both inside $S_\varepsilon$, and heights $\pm X_j$ from (E3).
By \eqref{eq:decay} and (E3), the horizontal sides carry $O(X_j^{-2}\log^2X_j)$; by $\Phi_h(1-s)=\Phi_h(s)$
($h$ even) and $\xi(s)=\xi(1-s)$ the left side reflects onto the right. On $\Rel s=\frac32$ write
$\xi'/\xi=\frac1s+\frac1{s-1}+\frac12\psi(\frac s2)-\frac12\log\pi+\zeta'/\zeta$. The two poles give
$\Phi_h(0)+\Phi_h(1)$. The archimedean part is holomorphic and $O(\log|s|)$ in $\Rel s>0$ (E4), and is
shifted to $\Rel s=\frac12$. For the prime part, $\sum\Lambda(n)n^{-3/2}<\infty$ (E1) licenses Fubini;
$\Phi_h(s)n^{-s}$ is holomorphic on $\frac12\le\Rel s\le\frac32$ with horizontals $O(X_j^{-2})$
uniformly in $n$, and (E5) yields $n^{-1/2}h(-\log n)$, doubled by evenness. The prime sum converges
since $|h(\pm\log n)|\le Cn^{-1-\varepsilon}$, and the zero sum since $|\Phi_h(\rho)|=O(|\gamma_\rho|^{-2})$
by \eqref{eq:decay} and (E3).
\end{proof}

Write $\langle h,\mathsf W\rangle$ for the prime sum minus the archimedean integral in
\eqref{eq:EF}, so that $\langle h,\mathsf W\rangle=\Phi_h(0)+\Phi_h(1)-\sum_\rho\Phi_h(\rho)$.

\begin{definition}\label{def:form}
For real $X\in\DW$ the function $X\star X^\vee$ is real and even; put $Q(X):=\langle X\star X^\vee,\mathsf W\rangle$ and
$\I(F,G):=\frac14[Q(F+G)-Q(F-G)]$.
\end{definition}

\begin{theorem}[comparison identity]\label{thm:comparison}
For all real $F,G\in\DW$, $\I$ is symmetric and bilinear and
\begin{equation}\label{eq:comparison}
\I(F,G)=\Phi_F(0)\Phi_G(1)+\Phi_F(1)\Phi_G(0)-\sum_\rho\Phi_F(\rho)\Phi_G(1-\rho),
\end{equation}
the zero sum absolutely convergent. Moreover $\I(F,G)=\langle h_{\rm ev},\mathsf W\rangle$ with
$h_{\rm ev}=\frac12(F\star G^\vee+F^\vee\star G)$.
\end{theorem}
\begin{proof}
$\Phi_{X\star X^\vee}(s)=\Phi_X(s)\Phi_X(1-s)$ by Lemma~\ref{lem:DW}(iii); substituting in \eqref{eq:EF}
gives $Q(X)=2\Phi_X(0)\Phi_X(1)-\sum_\rho\Phi_X(\rho)\Phi_X(1-\rho)$. Polarize term by term, using that the
zero multiset is invariant under $\rho\mapsto1-\rho$ (E2) and that the sum converges absolutely, so may
be re-indexed. The same substitution in $h_{\rm ev}$ gives the last statement.
\end{proof}

\begin{definition}\label{def:deg}
$\dg F:=\Phi_F(0)+\Phi_F(1)=\int F(t)\,2\cosh(t/2)\,dt$; $M(F):=(\Phi_F(1),\Phi_F(0))$;
$\PW:=\ker M$ (primitive classes); $F$ is effective, $F\ge0$, if $F(t)\ge0$ for all $t$; $\Ef$ is the
set of effective classes; $W(F):=\sum_\rho\Phi_F(\rho)\Phi_F(1-\rho)$, so $\I(F,F)=2\Phi_F(0)\Phi_F(1)-W(F)$.
\end{definition}

\begin{lemma}[$\Pi$ is a bijection onto the primitive class]\label{lem:Pi}
Let $I_T=(-T,T)$ and $\mathcal P_T=\{F\in C_c^\infty(I_T):\Phi_F(0)=\Phi_F(1)=0\}$. Then
$\Pi:C_c^\infty(I_T)\to\mathcal P_T$ is a linear bijection with inverse $F\mapsto e^{-|t|/2}\ast F$, and
$\Pi^{-1}$ preserves the convex hull of the support.
\end{lemma}
\begin{proof}
$\Pi$ maps into $\mathcal P_T$ since $\Phi_{\Pi G}=s(1-s)\Phi_G$, and is injective because $\Phi_G$ is entire.
$e^{-|t|/2}$ is the Green kernel of $\Pi$. For $F\in\mathcal P_T$ and $u=e^{-|t|/2}\ast F$, to the right of
$\supp F$ we have $u(t)=e^{-t/2}\Phi_F(1)=0$, and to the left $u(t)=e^{t/2}\Phi_F(0)=0$. So $u\in C_c^\infty$
with support in the hull of $\supp F$, and $\Pi u=F$.
\end{proof}

%======================================================================
\section{The P\'olya class and the hyperbolic polar plane}\label{sec:polya}

\begin{theorem}[the P\'olya class]\label{thm:K}
Let $K(t)=\sum_{n\ge1}(4\pi^2n^4e^{9t/2}-6\pi n^2e^{5t/2})e^{-\pi n^2e^{2t}}$.
\begin{enumerate}
\item[(i)] (see (X1)) $K$ is even and $\xi(\frac12+i\tau)=\int K(t)e^{i\tau t}dt$.
\item[(ii)] $K(t)>0$ for all $t$; $K(t)=4\pi^2e^{9|t|/2}e^{-\pi e^{2|t|}}(1+O(e^{-2|t|}))$, and
likewise for all derivatives with a polynomial factor in $e^{2|t|}$; hence $K$, all $K^{(j)}$ and all
$\tau_aK$ lie in $\DW$.
\item[(iii)] $\Phi_K=\xi$ on $\C$; $\Phi_K(0)=\Phi_K(1)=\frac12$, $M(K)=(\frac12,\frac12)$, $\dg K=1$.
\end{enumerate}
\end{theorem}
\begin{proof}
(ii) For $u\ge0$ each term equals $2\pi n^2e^{5u/2}e^{-\pi n^2e^{2u}}(2\pi n^2e^{2u}-3)>0$; positivity on
$\R$ follows by evenness. With $\omega(x)=\sum_{n\ge1}e^{-\pi n^2x}$ and $x=e^{2t}$,
$K(t)=4x^{9/4}\omega''(x)+6x^{5/4}\omega'(x)$. As $x\to\infty$ the $n=1$ term dominates. As $x\to0^+$,
the theta relation $\omega(x)=\frac12x^{-1/2}-\frac12+x^{-1/2}\omega(1/x)$ gives
$\omega'=-\frac14x^{-3/2}+O(x^{-5/2}e^{-\pi/x})$ and $\omega''=\frac38x^{-5/2}+O(x^{-9/2}e^{-\pi/x})$,
so the leading terms $\frac32x^{-1/4}-\frac32x^{-1/4}$ cancel and $K(t)=O(x^{-9/4}e^{-\pi/x})$. Each
$\partial_t$ costs a factor $O(e^{2|t|})$. (iii) $\Phi_K$ is entire by (ii) and agrees with $\xi$ on the
critical line by (i); $\xi(0)=\xi(1)=\frac12$ by (E2).
\end{proof}

\begin{theorem}[how $K$ pairs]\label{thm:Kpair}
For $F\in\DW$, $a\in\R$:
(i) $\I(F,K)=\frac12\dg F$; (ii) $\I(K,K)=\frac12$; (iii) $\I(F,\tau_aK)=\frac12(e^{a/2}\Phi_F(0)+e^{-a/2}\Phi_F(1))$;
(iv) $K\perp\PW$ and $K\ne0$; (v) $\I(F,2K)=\dg F$.
\end{theorem}
\begin{proof}
In \eqref{eq:comparison} with $G=K$ every zero term contains $\xi(1-\rho)=0$. (iii) uses
$\Phi_{\tau_aK}(s)=e^{(s-\frac12)a}\xi(s)$.
\end{proof}

\begin{remark}[on the name]\label{rem:name}
By analogy with Weil's proof one is tempted to call $K$ a canonical class. Section~\ref{sec:curve}
shows that under the curve dictionary the corresponding class is $\frac12(F_1+F_2)$, half the ample
ruling class, and not $K_{C\times C}=(2g-2)(F_1+F_2)$ (Theorem~\ref{thm:curveK}). We call $K$ the
\emph{P\'olya class}. Nothing below uses more than $\dg K>0$, the pairing of
Theorem~\ref{thm:Kpair}, and the affine symmetry $E\mapsto K-E$.
\end{remark}

\begin{lemma}\label{lem:Msurj}
For $a\ne0$, $M(K)=(\frac12,\frac12)$ and $M(\tau_aK)=\frac12(e^{a/2},e^{-a/2})$ are independent; $M$ is onto $\R^2$.
\end{lemma}

\begin{theorem}[the polar plane]\label{thm:split}
Fix $a>0$ and put $K_{\rm odd}=(2\sinh\frac a2)^{-1}(\tau_aK-\tau_{-a}K)$, so that
$\Phi_{K_{\rm odd}}(s)=\frac{\sinh((s-\frac12)a)}{\sinh(a/2)}\xi(s)$ and $M(K_{\rm odd})=(\frac12,-\frac12)$.
Let $\Lambda=\mathrm{span}(K,K_{\rm odd})$. Then
(i) $\PW=K^\perp\cap K_{\rm odd}^\perp$;
(ii) $\DW=\Lambda\oplus\PW$, $\I$-orthogonally;
(iii) $\I|_\Lambda$ has Gram matrix $\mathrm{diag}(\frac12,-\frac12)$ in the basis $(K,K_{\rm odd})$;
(iv) $\I(F,F)=-W(F)$ on $\PW$;
(v) $K^\perp=\{\dg F=0\}=\PW\oplus\R K_{\rm odd}$, and $\I\le0$ on $K^\perp$ iff $\I\le0$ on $\PW$.
\end{theorem}
\begin{proof}
$\Phi_{K_{\rm odd}}$ vanishes at every zero, so $\I(F,K_{\rm odd})=\frac12\Phi_F(0)-\frac12\Phi_F(1)$; with
$\I(F,K)=\frac12\dg F$ this gives (i). $M|_\Lambda$ is an isomorphism, which gives (ii). On $\Lambda$ only the
polar terms survive, giving (iii). (iv) is \eqref{eq:comparison} with $M(F)=0$, and (v) combines (ii)--(iv).
\end{proof}

\begin{proposition}[the polar block satisfies the Hodge-index inequality]\label{prop:slack}
Write $\mu(F)=\Phi_F(1)-\Phi_F(0)$. For every $F\in\DW$,
$\frac12(\dg F)^2-2\Phi_F(0)\Phi_F(1)=\frac12\mu(F)^2\ge0$. Hence $\I(F,F)\le\frac12(\dg F)^2$ holds at $F$
if and only if $W(F)\ge-\frac12\mu(F)^2$: no inequality between functions of $M(F)$ alone carries any
information about the zeros.
\end{proposition}

%======================================================================
\section{The radical}\label{sec:radical}

\begin{lemma}[separation lemma]\label{lem:sep}
Let $(c_\rho)$ be complex numbers indexed by the zeros with multiplicity, equal on copies of a
multiple zero, with $\sum|c_\rho|<\infty$, and let $\alpha,\beta\in\C$. If
$\sum_\rho c_\rho e^{(\rho-\frac12)t}=\alpha e^{t/2}+\beta e^{-t/2}$ for all real $t$, then $\alpha=\beta=0$ and all $c_\rho=0$.
\end{lemma}
\begin{proof}
Since $0<\Rel\rho<1$, $|e^{(\rho-\frac12)t}|\le e^{|t|/2}$ and both sides are continuous. Apply the
Laplace transform on $[0,\infty)$ at $z=s-\frac12$, $\Rel s>1$; termwise interchange is justified by
$\sum|c_\rho|<\infty$, and gives $H(s):=\sum_\rho\frac{c_\rho}{s-\rho}=\frac\alpha{s-1}+\frac\beta s$ for $\Rel s>1$.
$H$ is holomorphic off the discrete zero set, with at most a simple pole of residue $m_{\rho_0}c_{\rho_0}$
at a zero of multiplicity $m_{\rho_0}$. Both sides continue to $\C\setminus(Z(\xi)\cup\{0,1\})$, connected,
and $0,1\notin Z(\xi)$. Hence all residues at zeros vanish, $H\equiv0$, and $\alpha=\beta=0$.
\end{proof}

\begin{theorem}[the orthogonal of the primitive sector]\label{thm:perp}
For $G\in\DW$ the following are equivalent: (a) $\I(F,G)=0$ for all primitive $F\in C_c^\infty(\R)$;
(b) $\I(F,G)=0$ for all $F\in\PW$; (c) $\Phi_G(\rho)=0$ at every zero $\rho$.
\end{theorem}
\begin{proof}
(c)$\Rightarrow$(b) by \eqref{eq:comparison} and $\rho\mapsto1-\rho$; (b)$\Rightarrow$(a) is inclusion.
(a)$\Rightarrow$(c): with $c_\rho=\Phi_G(1-\rho)$, $\sum|c_\rho|<\infty$, and $\mathcal S(t)=\sum_\rho c_\rho e^{(\rho-\frac12)t}$ is
continuous with $\int F\mathcal S=0$ for primitive $F\in C_c^\infty$. By Lemma~\ref{lem:Pi},
$\langle\mathcal S,\Pi H\rangle=0$ for all $H\in C_c^\infty$, so $\Pi\mathcal S=0$ in $\mathcal D'(\R)$ and
$\mathcal S=\alpha e^{t/2}+\beta e^{-t/2}$. Lemma~\ref{lem:sep} gives $c_\rho=0$.
\end{proof}

\begin{remark}
Theorem~\ref{thm:perp}(c) asks for vanishing to order one at each zero. The inclusion
$\{G:\Phi_G|_{Z(\xi)}=0\}\supseteq\{G:\Phi_G=\xi e,\ e \text{ holomorphic on } S_\varepsilon\}$ is an
equality exactly when all zeros of $\xi$ are simple; we use only the left-hand side.
\end{remark}

\begin{theorem}[the radical]\label{thm:radical}
$\Nr:=\{N:\I(N,G)=0\ \forall G\in\DW\}=\{N:\Phi_N(0)=\Phi_N(1)=0,\ \Phi_N(\rho)=0\ \forall\rho\}=\PW\cap(\PW)^\perp$.
$\Nr$ is an ideal of $(\DW,\star)$ containing $\tau_a(\Pi K)$ for every $a$; on $\Nr$, $\dg=0$. Pairing to zero
against $C_c^\infty(\R)$ already forces membership in $\Nr$.
\end{theorem}
\begin{proof}
If $\I(N,G)=0$ for all $G\in C_c^\infty$, then as an identity of continuous functions
$\sum_\rho\Phi_N(\rho)e^{(\frac12-\rho)t}=\Phi_N(0)e^{t/2}+\Phi_N(1)e^{-t/2}$; re-index by $\rho\mapsto1-\rho$ and apply
Lemma~\ref{lem:sep}. The converse inclusion and the ideal property are immediate from \eqref{eq:comparison}
and $\Phi_{N\star g}=\Phi_N\Phi_g$.
\end{proof}

\begin{proposition}[$\Nr$ is not principal at $\Pi K$]\label{prop:notprincipal}
For $a\ne0$, $\tau_a(\Pi K)\in\Nr\setminus(\Pi K\star\DW)$.
\end{proposition}
\begin{proof}
If $\tau_a(\Pi K)=\Pi K\star g$, then on $\Rel s=\frac32$, where $s(1-s)\xi(s)\ne0$ by (E2),
$\Phi_g(\frac32+i\tau)=e^{(1+i\tau)a}$ has constant modulus $e^a$; but $g(t)e^t\in L^1$, so
$\Phi_g(\frac32+i\tau)\to0$ by Riemann--Lebesgue.
\end{proof}

\begin{theorem}[no compactly supported class is orthogonal to the primitive sector]\label{thm:compactrad}
$(\PW)^\perp\cap C_c^\infty(\R)=0$; in particular $\Nr\cap C_c^\infty(\R)=0$.
\end{theorem}
\begin{proof}
Let $N\in C_c^\infty$, $N\ne0$, with $\Phi_N(\rho)=0$ at every zero (Theorem~\ref{thm:perp}). By (E6),
$\Phi_N$ is entire of exponential type, so by (X2) it has $O(r)$ zeros in $|s-\frac12|\le r$. Every
\emph{distinct} zero of $\xi$ is a zero of $\Phi_N$, and by (X3) there are
$\ge\frac2{5\pi}r\log r(1+o(1))$ distinct zeros of $\xi$ in that disc. Contradiction.
\end{proof}

%======================================================================
\section{Degree, effective classes and bounded density}\label{sec:degree}

\begin{theorem}\label{thm:degree}
(i) If $E\ge0$, $E\ne0$, then $\dg E>0$. (ii) $\dg$ vanishes on $\Nr$ and on $\PW$; no nonzero primitive or
radical class is effective. (iii) $\dg=\I(\cdot,2K)$ and every $\I(\cdot,X)$ descend to $\DW/\Nr$.
\end{theorem}

\begin{theorem}[bounded density and ample classes]\label{thm:density}
For real $F\in\DW$ put $w_F(u)=\Phi_F(1)e^{-u/2}+\Phi_F(0)e^{u/2}-\sum_\rho e^{(\rho-\frac12)u}\Phi_F(1-\rho)$. Then:
(i) $\I(E,F)=\int E(u)w_F(u)\,du$ for every $E\in\DW$, absolutely convergent;
(ii) $C_F:=\sup_u|w_F(u)|/2\cosh(u/2)\le|\Phi_F(0)|+|\Phi_F(1)|+\sum_\rho|\Phi_F(1-\rho)|<\infty$, and
$|\I(E,F)|\le C_F\dg E$ for every effective $E$;
(iii) calling $H$ \emph{ample} if $\I(E,H)>0$ for all effective $E\ne0$, the class $2K$ is ample and so is
$F+2mK$ for every $m>C_F$.
\end{theorem}
\begin{proof}
(i) Expand \eqref{eq:comparison} with $\Phi_E(\rho)=\int Ee^{(\rho-\frac12)u}du$ and interchange, using
$|e^{(\rho-\frac12)u}|\le e^{|u|/2}$ and $\sum_\rho|\Phi_F(1-\rho)|<\infty$ (\eqref{eq:decay}, (E3)).
(ii) $e^{\pm u/2}\le2\cosh(u/2)$ and $|e^{(\rho-\frac12)u}|\le2\cosh(u/2)$, using only $0<\Rel\rho<1$.
(iii) $\I(E,2K)=\dg E>0$, and $\I(E,F+2mK)\ge(m-C_F)\dg E$.
\end{proof}

%======================================================================
\section{Affine duality and the Hodge-index transfer theorem}\label{sec:transfer}

The involution $\sigma_K(E):=K-E$ is affine, satisfies $\dg\sigma_K(E)=1-\dg E$, descends to $\DW/\Nr$, and
carries every primitive class to a class with $M=(\frac12,\frac12)\ne0$. Weil's duality $D\mapsto K_S-D$ is of
this type, a translation composed with negation, not an algebra involution.

\begin{theorem}[Hodge-index transfer]\label{thm:transfer}
Suppose $h^0:\DW\to\R_{\ge0}$ satisfies
\begin{enumerate}
\item[(a)] $h^0\ge0$;
\item[(b$'$)] if $h^0(E)>0$ there is $E'\in\DW$, $E'\ge0$, $E'\not\equiv0$, with $E-E'\in\Nr$;
\item[(c)] there is $\chi\in\R$ with $h^0(E)+h^0(K-E)\ge R(E):=\frac12\I(E,E)-\frac12\I(E,K)+\chi$ for every $E$.
\end{enumerate}
Then $\I(F,F)\le0$ for every real primitive $F\in C_c^\infty(\R)$.
\end{theorem}
\begin{proof}
Let $F$ be primitive with $q:=\I(F,F)>0$; then $\dg F=0$ and $\I(F,K)=0$.
\emph{Step 1.} Fix an integer $m>C_F$ and $H'=F+2mK$, ample by Theorem~\ref{thm:density}; $\I(F,H')=q$, $\I(K,H')=m$.
\emph{Step 2.} If $h^0(K-nF)>0$, (b$'$) gives $E'\ge0$, $E'\ne0$, $K-nF-E'\in\Nr$, so
$m-nq=\I(K-nF,H')=\I(E',H')>0$. Hence $h^0(K-nF)=0$ for $n>m/q$.
\emph{Step 3.} (c) at $E=nF$: $h^0(nF)\ge\frac12n^2q+\chi>0$ for $n$ large, since $h^0(K-nF)=0$.
\emph{Step 4.} (b$'$) gives $E'\ge0$, $E'\ne0$, with $nF-E'\in\Nr$, so $\dg E'=n\dg F=0$, contradicting Theorem~\ref{thm:degree}(i).
\end{proof}

\begin{remark}\label{rem:hartshorne}
This is \cite[V.1.9]{Hartshorne} with $H=2K$: ampleness of $2K$ and $F+2mK$ is Theorem~\ref{thm:density};
the vanishing of $h^0(K-nF)$ uses that $\Nr$ is the radical (Theorem~\ref{thm:radical}); Riemann--Roch is (c);
``effective of degree $0$ is $0$'' is Theorem~\ref{thm:degree}. The clause $E'\not\equiv0$ is used in Steps 2
and 4 and cannot be dropped. No sign $\chi\ge0$ is needed. The proof evaluates $h^0$ only at $nF$ and
$K-nF$. It remains valid if (b$'$) is weakened to: some positive multiple $NE$ is $\Nr$-equivalent to a
nonzero effective class (the factor $N>0$ passes through Steps 2 and 4).
\end{remark}

%======================================================================
\section{The Riemann--Roch functional}\label{sec:RR}

Let $R(E)=\frac12\I(E,E)-\frac14\dg E+\chi$ be the right side of (c) (by Theorem~\ref{thm:Kpair}(i)); by
\eqref{eq:comparison}, $R(E)=ab-\frac12W(E)-\frac14\dg E+\chi$ with $(a,b)=M(E)$.

\begin{theorem}[Serre symmetry]\label{thm:serre}
$R(K-E)=R(E)$ for every $E$ and every $\chi$; the proof uses only bilinearity and symmetry of $\I$.
\end{theorem}
\begin{proof}
$R(K-E)=\frac12[\I(K,K)-2\I(K,E)+\I(E,E)]-\frac12[\I(K,K)-\I(E,K)]+\chi=R(E)$.
\end{proof}

\begin{proposition}\label{prop:gated}
(i) $h^0=\max(R,0)$ satisfies (a) and (c) and fails (b$'$): $R(mK)=\frac{m^2}4-\frac m4+\chi\to\infty$ as
$m\to-\infty$ while $\dg(mK)<0$.
(ii) The gated candidate $h^0_+(E):=\max(R(E),0)\mathbf 1_{\{\dg E>0\}}$ satisfies (a) and (c) for every $\chi$,
vanishes in degree $\le0$, and has $h^0_+(mK)=\frac14m^2-\frac14m+\chi$ for large $m$.
(iii) If $\chi>0$, $h^0_+$ fails (b$'$).
\end{proposition}
\begin{proof}
(ii) $\dg E+\dg(K-E)=1$, so the two gates are not both off; use Theorem~\ref{thm:serre}.
(iii) For $c>0$, $\varepsilon>0$, $\beta=-6\varepsilon/(e^{c/2}-e^{-c/2})$, $\alpha=-2\varepsilon-\beta e^{c/2}$, the class
$E=\alpha K+\beta\tau_cK$ has $M(E)=(-\varepsilon,2\varepsilon)$, $W(E)=0$, $\dg E=\varepsilon$ and
$R(E)=-2\varepsilon^2-\frac\varepsilon4+\chi>0$ for small $\varepsilon$; but every nonzero effective class has
$\Phi(1)>0$, and $M$ is constant on $\Nr$-cosets.
\end{proof}

\begin{theorem}[the clause $(\mathrm b')^+$ suffices]\label{thm:bplus}
Let $\chi\in\R$ (for $\chi>0$ the hypothesis below is never satisfied, by
Proposition~\ref{prop:gated}(iii)) and suppose
\textup{(b$'$)$^+$}: every $E$ with $\dg E>0$ and $R(E)>0$ is $\Nr$-equivalent to a nonzero effective class.
Then the conclusion of Theorem~\ref{thm:transfer} holds. The same is true of the stable form, in which
some positive multiple of $E$ is required to be so.
\end{theorem}
\begin{proof}
First proof: $h^0_+$ satisfies (a), (c), and (b$'$)$^+$ is (b$'$) for it. Direct proof: let $F$ be primitive in
$C_c^\infty$ with $q=\I(F,F)>0$ and $E=rF+K$. Then $\dg E=1$ and $R(E)=\frac12r^2q+\chi>0$ for $r$ large;
an effective $E'\equiv E$ has $\dg E'=1$ and $\I(E',F)=rq$, while $|\I(E',F)|\le C_F$
(Theorem~\ref{thm:density}). Contradiction for $r>C_F/q$. In the stable form both sides scale by $N$.
\end{proof}

%======================================================================
\section{The frame defect theorem}\label{sec:defect}

\begin{definition}
$s(\gamma,\theta):=\sup_{u\in\R}|\cos(\gamma u+\theta)|/\cosh(u/2)\in(0,1]$.
\end{definition}

\begin{lemma}[the cosh inequality]\label{lem:cosh}
For $E'\in\DW$, $E'\ge0$, and all real $\gamma,\theta$:
$|\int E'(u)\cos(\gamma u+\theta)du|\le\frac12s(\gamma,\theta)\dg E'$. In particular, at a zero
$\rho=\frac12+i\gamma$, $|\Rel(e^{i\theta}\Phi_{E'}(\rho))|\le\frac12s(\gamma,\theta)\dg E'$.
\end{lemma}
\begin{proof}
$|\cos(\gamma u+\theta)|\le s\cosh(u/2)$ pointwise; integrate against $E'\ge0$.
\end{proof}

\begin{proposition}[the sharp constant at the first zero; computed in 40-digit arithmetic]\label{prop:s1}
$s(\gamma,\theta)<1$ for $\gamma\ne0$, $\theta\notin\pi\Z$, and
$s_1:=s(\gamma_1,\pi/2)=0.99846016271880910923\ldots$, attained at $u^*=0.110991563370513694\ldots$;
$1-s_1^2=3.0773034635\times10^{-3}$.
\end{proposition}
\begin{proof}
The numerator equals $1$ only on $(\pi\Z-\theta)/\gamma\not\ni0$, where the denominator exceeds $1$, and the
ratio tends to $0$ at infinity. For the value: outside $|u|\le0.4$ the ratio is $\le1/\cosh(0.2)<0.9804$;
inside, the critical equation $\gamma_1\cos(\gamma_1u)\cosh(u/2)=\frac12\sin(\gamma_1u)\sinh(u/2)$ is solved
by a 40-digit scan and a Newton step.
\end{proof}

\begin{lemma}[the test class; using (E2), (E3), (X4)]\label{lem:F0}
Let $\sigma=3$, $G=\sigma\sqrt{2\pi}/2=3.759942411946\ldots$ and $F_0(u)=e^{-u^2/18}\sin(\gamma_1u)$. Then:
(i) $F_0\in\DW$ is real and odd, so $\dg F_0=0$ exactly; $\Phi_{F_0}(0)=-\Phi_{F_0}(1)=\eta$, $|\eta|<10^{-389}$;
(ii) $\Phi_{F_0}(\rho_1)=iG(1+O(10^{-390}))$;
(iii) $L:=\sum_{\rho\ne\rho_1,\bar\rho_1}|\Phi_{F_0}(1-\rho)|\le5.5\times10^{-90}$;
(iv) $W_{F_0}:=-\I(F_0,F_0)=W(F_0)+2\eta^2$ satisfies $|W_{F_0}-2G^2|\le7\times10^{-89}$, where $2G^2=9\pi$;
(v) $C_{F_0}\le Gs_1+L+2|\eta|<Gs_1+10^{-89}$.
\end{lemma}
\begin{proof}
(i) $|\Phi_{F_0}(1)|\le\sigma\sqrt{2\pi}e^{\sigma^2(\frac14-\gamma_1^2)/2}<10^{-389}$. (ii)
$\Phi_{F_0}(\frac12+i\gamma)=\frac{\sigma\sqrt{2\pi}}{2i}(e^{-\sigma^2(\gamma+\gamma_1)^2/2}-e^{-\sigma^2(\gamma-\gamma_1)^2/2})$.
(iii) For $\rho=\beta+i\gamma$, $|\frac12-\beta|<\frac12$ by (E2), and
$|\Phi_{F_0}(1-\rho)|\le Ge^{9/8}(e^{-\sigma^2(\gamma-\gamma_1)^2/2}+e^{-\sigma^2(\gamma+\gamma_1)^2/2})$. By (X4) the other
zeros have $|\gamma|\ge21$, and there are at most $10\log T$ ordinates in $[T,T+1]$ for $T\ge21$; summing,
$L\le Ge^{9/8}\sum_{T\ge21}20\log T\,e^{-\frac92(T-\gamma_1)^2}\le5.5\times10^{-90}$.
(iv) By (X4) $\rho_1$ is simple; its pair contributes $2G^2(1+O(10^{-390}))$, the others at most
$\max_\rho|\Phi_{F_0}(\rho)|\cdot L\le11.6L$; and $\I(F_0,F_0)=-2\eta^2-W(F_0)$ by \eqref{eq:comparison}.
(v) $w_{F_0}(u)=-2G\sin(\gamma_1u)+r(u)$ with $|r(u)|\le(L+2|\eta|)\,2\cosh(u/2)$, and Proposition~\ref{prop:s1}.
\end{proof}

\begin{theorem}[a necessary condition and its failure]\label{thm:C41}
Let $F\in\DW$ be real with $\dg F=0$ and $W_F:=-\I(F,F)>0$, and $E_{t,d}=tF+dK$. Then $\dg E_{t,d}=d$,
$\I(E_{t,d},E_{t,d})=\frac{d^2}2-t^2W_F$, $\I(E_{t,d},F)=-tW_F$, and $R(E_{t,d})>0$ iff
$t^2W_F<\frac{d^2}2-\frac d2+2\chi$. If $E_{t,d}$ is $\Nr$-equivalent to an effective class, then
$tW_F\le C_Fd$. Consequently (b$'$)$^+$ forces $C_F\ge\sqrt{W_F/2}$. At $F=F_0$ this fails:
$C_{F_0}<Gs_1+10^{-89}<G(1-1.5\times10^{-3})<\sqrt{W_{F_0}/2}$.
\end{theorem}
\begin{proof}
Bilinearity with $\I(F,K)=0$ and $\I(K,K)=\frac12$; an effective $E'\equiv E_{t,d}$ has $\dg E'=d$ and
$\I(E',F)=-tW_F$, and Theorem~\ref{thm:density}(ii) applies. Let $d\to\infty$.
\end{proof}

\begin{theorem}[$(\mathrm b')^+$ is false; given (X4)]\label{thm:bfalse}
Put $s_{F}:=C_F\sqrt{2/W_F}$; by Lemma~\ref{lem:F0}, $s_{F_0}\le s_1^+:=s_1+10^{-89}$. Fix $\chi$ and $d$ with
$d^2(1-(s_1^+)^2)>d-4\chi$ (at $\chi=0$: $d>324.9598266\ldots$). Choose $t>0$ with
$s_{F_0}^2d^2/2<t^2W_{F_0}<\frac{d^2}2-\frac d2+2\chi$. Then $E=tF_0+dK$ has $\dg E=d$, $R(E)>0$, and neither
$E$ nor any positive multiple of $E$ is $\Nr$-equivalent to a nonzero effective class.
\end{theorem}
\begin{proof}
By Theorem~\ref{thm:C41}, equivalence would give $t^2W_{F_0}\le s_{F_0}^2d^2/2$; for multiples the factor cancels.
\end{proof}

\begin{remark}[an explicit instance]
At $\chi=0$, $d=400$, $t^2W_{F_0}=79780$ (so $t=53.11911\ldots$): $\dg E=400$, $\I(E,E)=220>\frac12\dg E$,
$R(E)=10$, and $|\Iml\Phi_E(\rho_1)|=tG=199.7248\ldots>\frac12s_1\dg E=199.6920\ldots$, whereas the
non-sharp bound $\sqrt{ab}=200$ is satisfied. The window is $((s_1^+)^2\cdot80000,79800)\supseteq(79753.82,79800)$.
\end{remark}

\begin{theorem}[frame defect; given (X4)]\label{thm:defect}
There is no $h^0:\DW\to\R_{\ge0}$ satisfying (a), (b$'$), (c) of Theorem~\ref{thm:transfer}, for any $\chi\in\R$.
\end{theorem}
\begin{proof}
\emph{Step 1.} If $h^0(E)>0$ then $\dg E=\dg E'>0$; so $h^0=0$ in degree $\le0$. \emph{Step 2.} If $\dg E>1$,
then $\dg(K-E)<0$, so $h^0(K-E)=0$, and (c) with Theorem~\ref{thm:serre} gives $h^0(E)\ge R(E)$.
\emph{Step 3.} The class $E$ of Theorem~\ref{thm:bfalse} with $d>1$ has $h^0(E)>0$, and (b$'$) supplies an
effective representative which does not exist.
\end{proof}

Theorem~\ref{thm:transfer} keeps its proof: it is an implication whose hypothesis is unsatisfiable. The
defect is localized: (c) with the Serre symmetry demands $h^0\ge R$ on the half-space $\dg>1$, while the
classes with effective representatives are strictly fewer than $\{R>0\}$ there; $\chi$ shifts $R$ by a
constant while the gap grows like $d^2(1-s_1^2)$.

\begin{lemma}[truncation; given (E2), (E3), (X4), (X5)]\label{lem:trunc}
Let $\mathcal N(g)=\int|g|\,2\cosh(u/2)$, $S_2=\sum_\rho|\rho-\frac12|^{-2}$, $\vartheta\in C_c^\infty$ even with
$0\le\vartheta\le1$, $\vartheta=1$ on $[-1,1]$, $\supp\vartheta\subseteq[-2,2]$, $A=1+\|\vartheta'\|_\infty+\|\vartheta''\|_\infty$, and
$F_R=F_0\vartheta(\cdot/R)$.
(i) $C_{F+F'}\le C_F+C_{F'}$, $C_F\le\mathcal N(F)+S_2\mathcal N(F'')$, and $|\I(E,F)|\le C_F\mathcal N(E)$ for all $E$.
(ii) $S_2\le0.0465$.
(iii) $F_R\in C_c^\infty$ is odd with $\dg F_R=0$ exactly.
(iv) For $R\ge6\sqrt{40+2\log A}$: $C_{F_R}\le Gs_1+10^{-6}+10^{-89}$, $|W_{F_R}-W_{F_0}|\le9.3\times10^{-5}$,
and $s_{F_R}\le s_1+3.6\times10^{-6}<0.998464$, so $(1-s_{F_R}^2)^{-1}<325.78$.
(v) With a fixed odd $B\in C_c^\infty$, $\Phi_B(0)=1$, and $R$ large, $F_R^\flat=F_R-\Phi_{F_R}(0)B$ is primitive,
odd, with $s_{F_R^\flat}\le0.998469$.
\end{lemma}
\begin{proof}
(i) Linearity of $w_F$; $|\Phi_F(1-\rho)|\le|\rho-\frac12|^{-2}\mathcal N(F'')$ from two integrations by parts at
$\Rel(1-\rho)$; Theorem~\ref{thm:density}(i). (ii) By (X5) $\sum\Rel(1/\rho)=0.0230957\ldots$; pairing $\rho$
with $1-\bar\rho$ gives $\sum(1+\gamma^2)^{-1}\le0.04620$, and $|\gamma|>14$ gives $S_2\le1.0052\cdot0.04620$.
(iii) Parity. (iv) With $D_R=F_0-F_R$, $|D_R''|\le Au^2e^{-u^2/18}\mathbf 1_{|u|\ge R}$, so
$\mathcal N(D_R)\le4e^{-R^2/36}$ and $C_{D_R}\le80ARe^{-R^2/36}\le10^{-6}$; then
$W_{F_R}-W_{F_0}=2\I(F_0,D_R)-\I(D_R,D_R)$ is bounded by (i). (v) Same estimates with $D'=D_R+\Phi_{F_R}(0)B$.
\end{proof}

\begin{theorem}[the defect on the small domain; given the inputs of Lemma~\ref{lem:trunc}]\label{cor:smalldomain}
Let $\mathcal D_c=C_c^\infty(\R)\oplus\R K$ and $\mathcal D_p=\mathcal P^\infty_c\oplus\R K$, where $\mathcal P^\infty_c$ is the
real primitive part of $C_c^\infty$. For $\mathcal D\in\{\mathcal D_c,\mathcal D_p\}$ there is no $h^0:\mathcal D\to\R_{\ge0}$
satisfying (b$'$) on $\mathcal D$ (with $E'\in\DW$) and (c) whenever $E,K-E\in\mathcal D$, for any $\chi$, also with
(b$'$) in stable form. More precisely, with $F=F_R$ or $F_R^\flat$ and $E_w=tF+dK$ as in
Theorem~\ref{thm:bfalse}, no $h^0\ge0$ satisfies (b$'$) at $E_w$ and $K-E_w$ and (c) at $E_w$.
\end{theorem}
\begin{proof}
$\dg(K-E_w)=1-d<0$ forces $h^0(K-E_w)=0$; (c) gives $h^0(E_w)\ge R(E_w)>0$; (b$'$) at $E_w$ is impossible
by Theorem~\ref{thm:bfalse} with $F$ in place of $F_0$, whose proof uses only $\dg F=0$, $W_F>0$ and $s_F<1$.
\end{proof}

\begin{remark}
(i) $\Nr\cap\mathcal D_c=0$ and the radical of $\I|_{\mathcal D_p}$ is $0$ (Theorem~\ref{thm:compactrad} and
Lemma~\ref{lem:sep}), so the passage $F_0\mapsto F_R$ is an analytic replacement, admissible because the
inequality is strict. (ii) The constant $s_*=\inf s_F$ over $\dg F=0$, $W_F>0$ may be computed over $C_c^\infty$,
by truncation, mollification and correction of the degree. (iii) Simplicity of $\rho_1$ is load-bearing: at
multiplicity $m$ one would have $C_{F_0}=mGs_1$ and $\sqrt{W_{F_0}/2}=G\sqrt m$, and $2s_1>\sqrt2$.
(iv) The same argument applies to any completed $L$-function with a known simple isolated zero on the line,
with $s_1$ replaced by the value of $s$ at that zero.
\end{remark}

\begin{theorem}[enlargement test; given (X4)]\label{thm:enlarge}
Let $\Ef'$ be any cone used in place of $\Ef$ in (b$'$) such that $\dg>0$ on $\Ef'\setminus\{0\}$ and
$C'_X:=\sup_{E'\in\Ef'}|\I(E',X)|/\dg E'<\infty$ for all $X$. Then the transfer hypothesis for $\Ef'$ is satisfiable
only if $C'_{F_0}\ge\sqrt{W_{F_0}/2}\ge G(1-10^{-88})$; in particular
$C'_{F_0}/C_{F_0}\ge(1-10^{-88})/(s_1+10^{-89})>1.00154$.
\end{theorem}
\begin{proof}
Theorems~\ref{thm:C41}--\ref{thm:defect} use only these two properties and bilinearity.
\end{proof}
Positive-type classes $G\star G^\vee$ fail the degree condition ($\dg=2\Phi_G(0)\Phi_G(1)=0$ for primitive $G$);
classes with $\Phi\ge0$ on the critical line fail both conditions; the cone on the translates $\tau_cK$ lies in
$\Ef$; closures cannot raise a support function. A cone defined by nonnegativity of the pairings
$\I(E,\tau_{\log n}g)$, $n\ge1$, for a fixed effective $g$, is not excluded by these remarks; whether it passes the test is not studied here.

%======================================================================
\section{A completed function with zeros off the line}\label{sec:xia}

Fix $0<a<\frac12$ and $\Xi_a(s)=\xi(s+a)\xi(s-a)$, with polar set $P=\{-a,a,1-a,1+a\}$ and
$K_a=(e^{at}K)\star(e^{-at}K)$, so that $\Phi_{K_a}=\Xi_a$, $K_a>0$, $K_a$ even.

\begin{proposition}[the frame for $\Xi_a$]\label{prop:xiaframe}
Let $\DW_a$ be Definition~\ref{def:DW} with decay exponent $1+a+\varepsilon$. Define $\I^{(a)}$ by the explicit
formula of $\Xi_a$, $\dg_a:=\I^{(a)}(\cdot,2K_a)=2\sum_{p\in P}\Xi_a(p)\Phi_F(p)$, primitivity by $\Phi_F(p)=0$ for $p\in P$.
Then Theorems~\ref{thm:comparison}, \ref{thm:Kpair}, \ref{thm:radical}, \ref{thm:degree}, \ref{thm:density},
\ref{thm:transfer} and \ref{thm:serre} hold for $\Xi_a$, with polar packet $\sum_{p\in P}\Phi_F(p)\Phi_G(1-p)$ and degree weight
$2\sum_p\Xi_a(p)e^{(p-\frac12)t}\asymp e^{(\frac12+a)|t|}$.
\end{proposition}
\begin{proof}
Three steps are not verbatim. (a) The contour abscissa: any $c\in(1+a,\frac32+a+\varepsilon)$ makes $\Phi_h$
holomorphic and $\sum\Lambda(n)(n^a+n^{-a})n^{-c}$ convergent; the four simple poles give the polar packet.
(b) The separation lemma holds on $\Rel s>1+a$, using $|\Rel\rho-\frac12|\le\frac12+a$, and needs $P\cap Z(\Xi_a)=\emptyset$,
which holds because $\Xi_a(\sigma)=\int K_ae^{(\sigma-\frac12)t}dt>0$ for real $\sigma$. (c) In
Theorem~\ref{thm:density}, $|e^{(\rho-\frac12)u}|\le e^{(\frac12+a)|u|}$ is dominated by the new degree weight. Theorem~\ref{thm:compactrad} is not claimed.
\end{proof}

\begin{lemma}\label{lem:PiP}
Let $c_1=\frac12+a$, $c_2=\frac12-a$ and $\Pi_P=(\partial_t^2-c_1^2)(\partial_t^2-c_2^2)$, of symbol $\pi_P(s)=\prod_{p\in P}(s-p)$.
Then $\Pi_P$ maps $C_c^\infty(I_T)$ bijectively onto the $\Xi_a$-primitive part of $C_c^\infty(I_T)$, with inverse
$F\mapsto G_{c_1}\ast G_{c_2}\ast F$, $G_c=\frac1{2c}e^{-c|t|}$, preserving support hulls and parity;
$\pi_P(1-s)=\pi_P(s)$; and no zero of $\Xi_a$ is real.
\end{lemma}
\begin{proof}
As in Lemma~\ref{lem:Pi}: $F_2=G_{c_2}\ast F$ vanishes outside the hull because $\Phi_F(a)=\Phi_F(1-a)=0$, and
$\Phi_{F_2}=\Phi_F/(c_2^2-(s-\frac12)^2)$ vanishes at $-a,1+a$ since $c_2^2-c_1^2=-2a\ne0$; repeat with $c_1$.
\end{proof}

\begin{proposition}\label{prop:xiaF}
If $\Xi_a$ has a zero $\rho_0$ with $\Rel\rho_0\ne\frac12$, there are nonzero real $\Xi_a$-primitive
$F_\pm\in C_c^\infty$, even and odd, with $\sum_{\rho\in Z(\Xi_a)}\Phi_{F_\pm}(\rho)\Phi_{F_\pm}(1-\rho)<0$, i.e.
$\I^{(a)}(F_\pm,F_\pm)>0$.
\end{proposition}
\begin{proof}
The proof is an adaptation of \cite[Thm.~6.2]{Paper45}, whose tests have the form
$q(i\,d/dt)\eta^{\star N}$ with $\eta$ an even bump; we indicate only the changes and do not write the
adapted construction out in full. The tests are taken in the form $\Pi_PG$ with $G$ of that type. The quartet
$\rho_0,1-\rho_0,\bar\rho_0,1-\bar\rho_0$ consists of four distinct points because $\Iml\rho_0\ne0$ (Lemma~\ref{lem:PiP}),
which replaces the use of $\zeta(\sigma)\ne0$ on $(0,1)$. Prescribing values of $\Phi_F=\pi_P\Phi_G$ at zeros is
equivalent to prescribing $\Phi_G$, since $\pi_P(\rho)\ne0$. The strip estimates are taken on
$|\Rel s-\frac12|\le\frac12+a$, valid for Paley--Wiener profiles on every strip. The tail acquires
$|\pi_P(\rho)|^2=O((1+|\rho|)^8)$, absorbed by raising the polynomial degree of the construction by $4$;
the zero count of $\Xi_a$ is $2N(T)+O(1)$.
\end{proof}

\begin{theorem}[no $h^0$ for $\Xi_a$; using (X6)]\label{thm:xia-noh}
For $\Xi_a$ no $h^0$ satisfies (a), (b$'$), (c). Moreover (b$'$)$^+$ fails for $\Xi_a$ at every $\chi$, with witness
$E=rF_++K_a$, $r$ large.
\end{theorem}
\begin{proof}
By (X6), $\Xi_a$ has infinitely many zeros with real part $\frac12\pm a$; by Proposition~\ref{prop:xiaF} there is a
primitive $F_+$ with $q=\I^{(a)}(F_+,F_+)>0$, contradicting the conclusion of Theorem~\ref{thm:transfer} for
$\Xi_a$. (This part rests on Proposition~\ref{prop:xiaF}, whose proof is given only as an
adaptation of \cite[Thm.~6.2]{Paper45}.) For the witness, $\dg_aE=\dg_aK_a>0$, $R_a(E)=\frac12r^2q+\frac12\I^{(a)}(K_a,K_a)-\frac14\dg_aK_a+\chi>0$, and
an effective representative $E'$ would give $rq=|\I^{(a)}(E',F_+)|\le C^{(a)}_{F_+}\dg_aK_a$.
\end{proof}

%======================================================================
\section{The function-field model}\label{sec:curve}

Let $C/\F_q$ be smooth, projective, geometrically irreducible of genus $g$, with
$Z_C(T)=P(T)/((1-T)(1-qT))$, $P(T)=\prod_{i=1}^{2g}(1-\alpha_iT)$, $N_m=q^m+1-\sum_i\alpha_i^m$, $|\alpha_i|=\sqrt q$ (X7).
On $S=C\times C$ with fibres $F_1=\{P\}\times C$, $F_2=C\times\{P\}$, diagonal $\Delta$ and graphs $\Gamma_f$:
$F_i^2=0$, $F_1\cdot F_2=1$, $\Delta^2=2-2g$, $\Gamma_f\cdot F_1=1$, $\Gamma_f\cdot F_2=\deg f$,
$K_S=(2g-2)(F_1+F_2)$, $\chi(\mathcal O_S)=(1-g)^2$ (X8).

\subsection{The discrete frame}
\begin{definition}\label{def:curveframe}
With $z=q^{s-1/2}$ and $u_i=\alpha_i/\sqrt q$, put $\Lambda_C(z)=z^{-g}\prod_i(z-u_i)$ and $\xi_C(s)=\Lambda_C(q^{s-1/2})$.
Let $\DW_C=\{E:\Z\to\R:\ E(n)=O(q^{-(1+\varepsilon)|n|})\}$, $\Phi_E(s)=\sum_nE(n)q^{(s-\frac12)n}$, and
$\dg E=\Phi_E(0)+\Phi_E(1)=\sum_nE(n)\,2\cosh(\frac n2\log q)$.
\end{definition}
The idele class group of $\F_q(C)$ has modulus group $q^\Z$, so the modulus line becomes the lattice $(\log q)\Z$.

\begin{lemma}\label{lem:xiC}
$\xi_C(s)=\xi_C(1-s)$; its zeros are the $\rho_i$ with $q^{\rho_i}=\alpha_i$, on $\Rel s=\frac12$; and
$\xi_C(0)=\xi_C(1)=hq^{-g/2}$ with $h=P(1)$.
\end{lemma}
\begin{proof}
$\prod u_i=1$ and the multiset $\{u_i\}$ is stable under conjugation, so $\Lambda_C(1/z)=\Lambda_C(z)$. Using $\alpha\mapsto q/\alpha$,
$\prod(q-\alpha_i)=q^gP(1)$, whence $\Lambda_C(q^{1/2})=q^{-g/2}h$.
\end{proof}

\begin{theorem}[explicit formula and dictionary; from (X7), (X8)]\label{thm:dict}
Let $\kappa(n)=q^{-|n|/2}N_{|n|}$ with $N_0=2-2g$. For $E\in\DW_C$,
$\Phi_E(0)+\Phi_E(1)-\sum_i\Phi_E(\rho_i)=\sum_nE(n)\kappa(n)$. Put
$\I^C(E,F)=\Phi_E(0)\Phi_F(1)+\Phi_E(1)\Phi_F(0)-\sum_i\Phi_E(\rho_i)\Phi_F(1-\rho_i)$ and
$D_E=\sum_nE(n)\gamma_n$, $\gamma_n=q^{-|n|/2}\Gamma_{q^n}$ $(n\ge0)$, $\gamma_{-n}=q^{-n/2}\Gamma_{q^n}^t$. Then
(i) $\gamma_a\cdot\gamma_b=\kappa(a-b)$; (ii) $\I^C(E,F)=D_E\cdot D_F$; (iii) $D_E\cdot F_1=\Phi_E(0)$,
$D_E\cdot F_2=\Phi_E(1)$, $\dg E=D_E\cdot(F_1+F_2)$; (iv) $E$ is primitive iff $D_E\perp F_1,F_2$.
\end{theorem}
\begin{proof}
$\sum_i\Phi_E(\rho_i)=\sum_nE(n)q^{-|n|/2}\sum_i\alpha_i^{|n|}$ and $\Phi_E(0)+\Phi_E(1)=\sum_nE(n)(q^{|n|/2}+q^{-|n|/2})$;
subtract and use $N_m=q^m+1-\sum\alpha_i^m$, also at $m=0$. (i) is the Lefschetz relations; (ii) expands
$D_E\cdot D_F=\sum_m(E\star F^\vee)(m)\kappa(m)$ and applies the explicit formula to $E\star F^\vee$; (iii) is
$\Gamma_{q^n}\cdot F_1=1$, $\Gamma_{q^n}\cdot F_2=q^n$.
\end{proof}
So the half-density weight $q^{-|n|/2}$ on the graphs makes the Gram Toeplitz and the sign $+$; the target
inequality $\I^C(F,F)\le0$ on primitives is the Hodge index theorem on $\{F_1,F_2\}^\perp$ with $H=F_1+F_2$.

\begin{theorem}[radical and isometry]\label{thm:curverad}
If the $\alpha_i$ are distinct, the radical of $\I^C$ is the common kernel of $\Phi(0),\Phi(1),\Phi(\rho_i)$,
$\dim\DW_C/\Nr^C=2+2g$, and $E\mapsto D_E$ is an isometric embedding into $V_C=\sum_n\R\gamma_n\subseteq NS(S)_\R$,
of signature $(1,1+2g)$.
\end{theorem}
\begin{proof}
Pairing $N$ with each $\delta_m$ gives $\Phi_N(0)q^{m/2}+\Phi_N(1)q^{-m/2}=\sum_i\Phi_N(\rho_i)\bar u_i^m$, an identity
between distinct characters of $\Z$; the same independence gives surjectivity of the evaluation map from finitely supported real $E$ onto $\{(a,b,(z_i)):a,b\in\R,\ z_j=\bar z_i\text{ when }u_j=\bar u_i\}$ (take the real part of a complex preimage).
\end{proof}

For $g=1$ write $(a_E,b_E,x_E,y_E)=(\Phi_E(1),\Phi_E(0),\Rel\Phi_E(\rho_1),\Iml\Phi_E(\rho_1))$; then
$\I^C(E,F)=a_Eb_F+b_Ea_F-2(x_Ex_F+y_Ey_F)$. For ordinary $E/\F_q$, $\mathrm{rk}\,NS(E\times E)=4$ (X7), so
$V_C=NS(E\times E)_\R$ and the rulings $F_1,F_2$ are classes of the model.

\subsection{The defect on the curve}
Let $\Ef_C=\{E\in\DW_C:E\ge0,\ E\ne0\}$ and $\Kc$ the class of Theorem~\ref{thm:curveK}(iii).

\begin{theorem}\label{thm:curveE}
(i) $\dg>0$ on $\Ef_C$, $\dg=0$ on $\Nr^C$, and $\I^C(E,2\Kc)=\dg E$. (ii) With
$w_F(n)=\Phi_F(1)q^{-n/2}+\Phi_F(0)q^{n/2}-\sum_iq^{(\rho_i-\frac12)n}\Phi_F(1-\rho_i)$, $|\I^C(E,F)|\le C_F\dg E$ on $\Ef_C$
with $C_F\le|\Phi_F(0)|+|\Phi_F(1)|+\sum_i|\Phi_F(1-\rho_i)|$. (iii) $\Ef_C$ is generated by the extreme rays $\delta_n$, and
$D_{\delta_n}=\gamma_n$: the frame's effective cone is the cone on $q^{-n/2}\Gamma_{q^n}$ and $q^{-n/2}\Gamma_{q^n}^t$, $n\ge0$.
\end{theorem}

\begin{lemma}[cosh inequality on the lattice]\label{lem:curvecosh}
With $s_C(\gamma,\theta)=\sup_{n\in\Z}|\cos(\gamma n\log q+\theta)|/\cosh(\frac n2\log q)$, every $E'\in\Ef_C$ satisfies
$|\Rel(e^{i\theta}\Phi_{E'}(\frac12+i\gamma))|\le\frac12s_C(\gamma,\theta)\dg E'$. Moreover $s_C<1$ iff $\theta\notin\pi\Z$.
\end{lemma}

\begin{proposition}[the sharp constant; exact]\label{prop:sC}
Let $g=1$, $a=q+1-N_1$, $\theta_u=\gamma_1\log q$. At $\theta=\pi/2$, $s_C^2=(4q-a^2)/(q+1)^2$ when the $n=\pm1$ entry
dominates. Minimizing over $\theta$ by balancing the entries $n=0,1$ gives
$\tan\theta^*=(q+1+a)/\sqrt{4q-a^2}$ and $(s_C^*)^2=(4q-a^2)/((4q-a^2)+(q+1+a)^2)$, provided no entry with
$|n|\ge2$ exceeds it. For the curve $E_0/\F_2:y^2+xy=x^3+1$ ($N_1=4$, $a=-1$, $h=4$, ordinary): $s_C(\gamma_1,\pi/2)=\sqrt7/3$ and
$s_C^*=\sqrt{7/11}=0.797724\ldots$, with $(1-(s_C^*)^2)^{-1}=\frac{11}4$ exactly; the supremum at $\theta^*$ is a triple tie at
$n=-2,0,1$, and all other entries are $\le0.620453$.
\end{proposition}
\begin{proof}
$\cos\theta_u=a/2\sqrt q$, $\cosh(\frac12\log q)=(q+1)/2\sqrt q$; solve $|\cos\theta|\cosh(\frac12\log q)=|\cos(\theta_u+\theta)|$.
For the curve, $\cos\theta^*=\sqrt{7/11}$, $\cos2\theta_u=-\frac34$, $\sin2\theta_u=-\frac{\sqrt7}4$, and the $n=-2$ entry is
$\frac{5\sqrt7}{4\sqrt{11}}\cdot\frac45=\sqrt{7/11}$; entries with $|n|\ge5$ are $\le1/\cosh(\frac52\log2)<0.18$.
\end{proof}

\begin{theorem}[the witness and the defect on the curve]\label{thm:curvedefect}
Assume the $\alpha_i$ distinct. Fix $\chi$, a direction $\theta$ with $s_C=s_C(\gamma_1,\theta)$, and $d>1$ with $d^2(1-s_C^2)>d-4\chi$ (such $d$ exist for every $\chi$ when $\theta\notin\pi\Z$). Let $\rho_1=\frac12+i\gamma_1$ be a zero of $\xi_C$; then $u_1=q^{i\gamma_1}$ is not real (the real $u_i$ would be distinct elements of $\{\pm1\}$, even in number, hence $\{1,-1\}$, and $\prod_iu_i$ would be $-1$), so $\bar\rho_1$ is a different zero modulo $2\pi i/\log q$; let $F_\theta$ be the
class with $\Phi(0)=\Phi(1)=0$, $\Phi(\rho_1)=e^{-i\theta}$, $\Phi(\bar\rho_1)=e^{i\theta}$ and $\Phi(\rho_i)=0$ at the other zeros, which exists by Theorem~\ref{thm:curverad} (for $g=1$: coordinates $(0,0,\cos\theta,-\sin\theta)$), so $\Rel(e^{i\theta}\Phi_{F_\theta}(\rho_1))=1$ and $-\I^C(F_\theta,F_\theta)=2$, and
$E=tF_\theta+d\Kc$ with $\frac12s_Cd<t<\frac12\sqrt{d(d-1)+4\chi}$. Then $\dg E=d$, $R(E)=\frac{d^2}4-t^2-\frac d4+\chi>0$, and
$E$ has no representative in $\Ef_C$ modulo $\Nr^C$. Consequently no $h^0:\DW_C\to\R_{\ge0}$ satisfies (a), (b$'$), (c)
with $K=\Kc$, cone $\Ef_C$ and radical $\Nr^C$, for any $\chi$. For the curve $E_0$ and $\chi=0$, $d=3$ suffices at $\theta^*$
(window $t\in(1.19659,1.22474)$), and $d=5$ at $\theta=\pi/2$.
\end{theorem}
\begin{proof}
Since $\Phi_{\Kc}$ vanishes at every $\rho_i$ and $\Phi(1-\rho)=\Phi(\bar\rho)$ on the critical line, $E$ has $\Phi_E(0)=\Phi_E(1)=\frac d2$, $\Phi_E(\rho_1)=te^{-i\theta}$, $\Phi_E(\bar\rho_1)=te^{i\theta}$ and $\Phi_E(\rho_i)=0$ otherwise (for $g=1$: coordinates $(\frac d2,\frac d2,t\cos\theta,-t\sin\theta)$), so $\I^C(E,E)=\frac{d^2}2-2t^2$ by Theorem~\ref{thm:dict}. A representative
$E'\in\Ef_C$ would have $\dg E'=d$ and $|\Rel(e^{i\theta}\Phi_{E'}(\rho_1))|=t>\frac12s_Cd$, against
Lemma~\ref{lem:curvecosh}. The non-existence of $h^0$ follows as in Theorem~\ref{thm:defect}.
\end{proof}

\begin{theorem}[the geometric $h^0$; (X8)]\label{thm:geomh0}
$h^0_C(D)=\dim H^0(S,\mathcal O(D))$ satisfies (a), (b$'$) with $\mathrm{Eff}(S)$, and (c) with $K_S$ and
$\chi=(1-g)^2$; Weil's proof is Theorem~\ref{thm:transfer} run on $(h^0_C,\mathrm{Eff}(S),K_S)$.
\end{theorem}

Of the three differences between Theorems~\ref{thm:curvedefect} and \ref{thm:geomh0} only one matters. The
class $K$ enters only through $\dg K>0$ and the symmetry of $R$; $\chi$ is quantified over all reals. The
cone is the whole gap: the defect is a property of $\Ef_C$, present where $h^0$ exists and the Hodge index
theorem holds.

\subsection{The cones}
\begin{theorem}[$g=1$]\label{thm:cones}
Let $\sigma(\mathcal C)=\sup\{|y_D|/\dg D:D\in\mathcal C,\dg D>0\}$. Then:
(i) $\sigma(\Ef_C)=\frac12s_C(\gamma_1,\pi/2)$; when the $n=\pm1$ entry dominates, as it does for the curve $E_0$, it is attained at $\gamma_{\pm1}$ ($=\sqrt7/6$ for $E_0$);
(ii) the positive cone $\{D^2\ge0,\ D\cdot H>0\}$ has $\sigma=\frac12$;
(iii) $\overline{\mathrm{Eff}}(E\times E)=\overline{\mathrm{Pos}}$, so $\sigma(\overline{\mathrm{Eff}})=\frac12$;
(iv) in each direction $\theta$ the ratio of support functions of $\Ef_C$ and $\overline{\mathrm{Eff}}$ is $s_C(\gamma_1,\theta)$;
(v) $F_1\in\mathrm{Eff}\setminus\Ef_C$, since every nonzero class of $\Ef_C$ has $\Phi(0)>0$ while $F_1$ has $\Phi(0)=F_1\cdot F_1=0$.
Adjoining $F_1,F_2$ to $\Ef_C$ does not change $\sigma$.
\end{theorem}
\begin{proof}
(i) A linear functional attains its supremum on the extreme rays, which give $|\sin(n\theta_u)|/2\cosh(\frac n2\log q)$.
(ii) Maximize $|y|/(a+b)$ under $ab\ge x^2+y^2$: AM--GM. (iii) Big classes are effective up to multiples (X8);
on the abelian surface $E\times E$ effective divisors are nef, so $\overline{\mathrm{Eff}}\subseteq\overline{\mathrm{Pos}}$.
\end{proof}
For $g\ge2$ the upper bound in (iii) is open: the surface has negative curves ($\Delta^2=2-2g$).

\begin{theorem}[the Frobenius-nef cone]\label{thm:B4C}
Let $B4_C=\{E:\I^C(E,\delta_n)\ge0\ \forall n\}$. Then $B4_C$ is the dual of $\Ef_C$, contains
$\mathrm{Nef}(S)\cap V_C$, and $\Ef_C\subseteq B4_C$ iff $\kappa\ge0$ iff $g\le1$. The Gram of the generators of $\Ef_C$ is
$q^{-|a-b|/2}N_{|a-b|}$, the data of the Euler product $Z_C(T)=\exp\sum N_mT^m/m$.
\end{theorem}

\subsection{The P\'olya class of the curve is half the ample class}
\begin{theorem}\label{thm:curveK}
Let $K_C(n)=k_n$ with $\Lambda_C(z)=\sum_{|n|\le g}k_nz^n$. Then (i) $\Phi_{K_C}=\xi_C$; (ii)
$\Phi_{K_C}(0)=\Phi_{K_C}(1)=hq^{-g/2}$; (iii) $\Kc=\frac{q^{g/2}}{2h}K_C$ has $\dg\Kc=1$, $\I^C(F,\Kc)=\frac12\dg F$,
$\I^C(\Kc,\Kc)=\frac12$; (iv) $D_{\Kc}\equiv\frac12(F_1+F_2)$; (v) $K_S^2=2(2g-2)^2\ne\frac12$, so $\Kc$ is not the
canonical class of $S$.
\end{theorem}
\begin{proof}
(iii) The zero terms vanish by (i). (iv) Both classes pair as $(\frac12,\frac12,0)$ with $F_1$, $F_2$ and the zero
block, which determine a class modulo the radical.
\end{proof}

\begin{remark}
$K>0$ has no general analogue: for $g=1$, $K_C=\delta_1+k_0\delta_0+\delta_{-1}$ with $k_0=-a/\sqrt q$, so
$K_C\ge0$ iff $N_1\ge q+1$. This holds for the curve $E_0$ ($N_1=4$, $k_0=1/\sqrt2$) and fails, for
example, for every elliptic curve over $\F_2$ with $N_1\le2$. The frame uses only $\dg\Kc=1$.
\end{remark}

\subsection{The missing sums}
For an ordinary elliptic curve with Frobenius $\pi$, $NS(E\times E)_\R$ is spanned by $F_1,F_2$ and the graphs
$\Gamma_\alpha$, $\alpha\in\Z[\pi]$, with $\Gamma_\alpha\cdot\Gamma_\beta=\deg(\alpha-\beta)$ (X7). Since $N_m=\deg(1-\pi^m)$,
the Gram of Theorem~\ref{thm:dict}(i) is the case $\alpha=\pi^a$, $\beta=\pi^b$: the frame's cone is the cone on the
multiplicative semigroup $\{\pi^n\}$ and its transposes. The missing effective classes are the graphs of sums
$m+n\pi$ and the fibres. In these terms Theorem~\ref{thm:cones}(iii) is the positive semidefiniteness of
$\deg(m+n\pi)=m^2+amn+qn^2$, i.e. $a^2\le4q$; the form is already determined by the Gram ($a$ is read from
$\kappa(1)$), and its positivity comes from $\deg$ being a count, $\#\ker(m+n\pi)$. The diagonal entry
$\kappa(0)=2-2g$ tends to $-\infty$ with the genus; a transport of the sums to the arithmetic frame would have
to be a difference-only form or carry an explicit counterterm. Further mismatches between the model and the
arithmetic frame: the model has no archimedean place (its archimedean term is the integer $2-2g$); it has no
smoothness condition, so point masses are classes and are the extreme rays of the cone; periodicity makes
the quotient by the radical finite-dimensional; and on the lattice the optimal direction is not $\pi/2$,
because $n=0$ is a sample point.

%======================================================================
\section{What the frame lacks}\label{sec:open}

In geometry the effective cone fills the positive cone. In $\DW$ effectivity is pointwise nonnegativity of a
single function on the modulus line, so the image of $\Ef$ under $(\dg,(\Phi(\rho))_\rho)$ is a moment cone, and
the cosh inequality measures how much thinner it is: the degree weight grows while the zero weight
oscillates. The frame retains the intersection theory of $C\times C$ and discards its linear systems.
The model shows this is a property of the cone and not of the zeros. Two questions are left open:
\begin{enumerate}
\item Is there a cone $\Ef'\supseteq\Ef$ (modulo $\Nr$) with $\dg>0$ off $\Nr$, finite support functions, and
$C'_{F}\ge\sqrt{W_F/2}$ for all $F$ of degree $0$ with $W_F>0$, defined without reference to the zeros?
Theorem~\ref{thm:enlarge} is the test any candidate must pass. In the model, the graphs of sums supply it.
\item What is $\sigma(\overline{\mathrm{Eff}})$ for curves of genus $g\ge2$, and does a curve model with an
archimedean place exist in which the same comparison can be made?
\end{enumerate}

\appendix
\section{Numerical constants}
The constants of Propositions~\ref{prop:s1} and \ref{prop:sC}, of Lemma~\ref{lem:F0} and of the instance
following Theorem~\ref{thm:bfalse} are reproduced by the script \texttt{scripts/constants.py}
distributed with this paper (\texttt{mpmath}, 40 digits).

\begin{thebibliography}{99}
\bibitem{Barner1981} K.~Barner, \emph{On A.~Weil's explicit formula}, J.~reine angew.\ Math.\ \textbf{323} (1981), 139--152.
\bibitem{Conrey1989} J.~B.~Conrey, \emph{More than two fifths of the zeros of the Riemann zeta function are on the critical line},
J.~reine angew.\ Math.\ \textbf{399} (1989), 1--26.
\bibitem{Davenport2000} H.~Davenport, \emph{Multiplicative Number Theory}, 3rd ed., Grad.\ Texts in Math.\ \textbf{74}, Springer, New York, 2000.
\bibitem{Hardy1914} G.~H.~Hardy, \emph{Sur les z\'eros de la fonction $\zeta(s)$ de Riemann}, C.~R.~Acad.\ Sci.\ Paris \textbf{158} (1914), 1012--1014.
\bibitem{Hartshorne} R.~Hartshorne, \emph{Algebraic Geometry}, Grad.\ Texts in Math.\ \textbf{52}, Springer, New York, 1977.
\bibitem{HeathBrown1979} D.~R.~Heath-Brown, \emph{Simple zeros of the Riemann zeta-function on the critical line}, Bull.\ London Math.\ Soc.\ \textbf{11} (1979), 17--18.
\bibitem{Lazarsfeld} R.~Lazarsfeld, \emph{Positivity in Algebraic Geometry I}, Ergeb.\ Math.\ Grenzgeb.\ \textbf{48}, Springer, Berlin, 2004.
\bibitem{Levinson1974} N.~Levinson, \emph{More than one third of zeros of Riemann's zeta-function are on $\sigma=1/2$}, Adv.\ Math.\ \textbf{13} (1974), 383--436.
\bibitem{Paper45} D.~A.~Trejo Pizzo, \emph{An explicit formula for the trace on semilocal idele class groups}, preprint, 2026.
\bibitem{Paper46} D.~A.~Trejo Pizzo, \emph{The sign of the radial intersection form on the idele class group}, preprint, 2026 (the preceding paper of this series).
\bibitem{PlattTrudgian2021} D.~J.~Platt and T.~S.~Trudgian, \emph{The Riemann hypothesis is true up to $3\cdot10^{12}$}, Bull.\ London Math.\ Soc.\ \textbf{53} (2021), 792--797.
\bibitem{Polya1926} G.~P\'olya, \emph{Bemerkung \"uber die Integraldarstellung der Riemannschen $\xi$-Funktion}, Acta Math.\ \textbf{48} (1926), 305--317.
\bibitem{Tate1966} J.~Tate, \emph{Endomorphisms of abelian varieties over finite fields}, Invent.\ Math.\ \textbf{2} (1966), 134--144.
\bibitem{Titchmarsh1986} E.~C.~Titchmarsh, \emph{The Theory of the Riemann Zeta-Function}, 2nd ed., revised by D.~R.~Heath-Brown, Oxford Univ.\ Press, 1986.
\bibitem{Trudgian2014} T.~S.~Trudgian, \emph{An improved upper bound for the argument of the Riemann zeta-function on the critical line II}, J.~Number Theory \textbf{134} (2014), 280--292.
\bibitem{Weil48} A.~Weil, \emph{Sur les courbes alg\'ebriques et les vari\'et\'es qui s'en d\'eduisent}, Hermann, Paris, 1948.
\bibitem{Weil52} A.~Weil, \emph{Sur les ``formules explicites'' de la th\'eorie des nombres premiers}, Comm.~S\'em.~Math.~Univ.~Lund (1952), 252--265.
\end{thebibliography}
\end{document}
